% Historical proof from Overleaf 5bfd38b4cc0b17c77f37ca07c08603ff6b06f309.
% Not in the current manuscript; uses its macros and references.
% Original certificates/verify_spectral.py is now verify_spectral.py here.
\begin{proposition}\label{prop:spectral-redundancy}
There are admissible rules with the same scale, edge count, thermal multiplier and mass Perron root, but different full three-state spectra.
\end{proposition}
\begin{proof}
Use the $19$-based allocation at depth $424$. For the four words with common prefix $0^{337}1^{83}$ and tails $1100,1010,1001,0110$, change their allocation counts by
\[
 -17413779,\qquad169010230,\qquad-155094721,\qquad3498270,
\]
respectively. All four words have $339$ zeroes and begin with $00$. They therefore avoid every correction packet, whose first pair is $01$ or $10$. Their initial counts are the zero-first-letter count at $z=339$, possibly increased by one. The supplied certificate satisfies
\[
 \min\{\text{that count},\ 2^{339}-\text{that count}-1\}
 \ge169010230.
\]
Thus both the displayed changes and their negatives preserve every capacity bound. This inequality and the small matrix identities below are checked in \path{certificates/verify_spectral.py}.

The changes sum to zero and have the same zero count, so they preserve the edge count and thermal multiplier. They do not affect the all-zero word and hence preserve the distance. The weighted sum of the four tail matrices, multiplied on the right by $\mathbf D$, is
\[
 \begin{pmatrix}
 304626366720&-728454355200\\
 119886741000&-286685685000
 \end{pmatrix}
 =\begin{pmatrix}13244624640\\5212467000\end{pmatrix}(23,-55).
\]
This matrix annihilates $(55,23)^T$, so left multiplication by the prefix matrix preserves the positive mass eigenvector and its eigenvalue. Since the changed allocation still defines an admissible graph, that eigenvalue remains the Perron root.

The trace change of the two-dimensional block is
\[
 \frac{(23,-55)\mathbf K^{337}\mathbf J^{83}
       (13244624640,5212467000)^T}{16^{425}}<0.
\]
Indeed, $(23,-55)\mathbf K^5=(-4444644,-17343936)$, and all entries of $\mathbf K$, $\mathbf J$ and the right-hand column are positive. The scalar thermal eigenvalue is unchanged, so the trace change of the full three-state matrix equals that of its two-dimensional block. The original rule and its two opposite perturbations have three different traces, although all four growth multipliers agree. Their spectra are therefore different.
\end{proof}
